\documentclass[11pt,a4paper]{article}
\usepackage{amsmath,amssymb}
\usepackage{listings}
\usepackage{hyperref}
\usepackage[margin=1in]{geometry}
\usepackage{fancyhdr}

\pagestyle{fancy}
\fancyhf{}
\fancyhead[L]{Module 19: Scientific Computing}
\fancyhead[R]{C Programming Course}
\fancyfoot[C]{\thepage}

\lstset{
  language=C,
  basicstyle=\ttfamily\small,
  keywordstyle=\bfseries,
  commentstyle=\itshape,
  numbers=left,
  numberstyle=\tiny,
  frame=single,
  breaklines=true,
  tabsize=4
}

\title{Module 19: Scientific Computing}
\author{C Programming Course}
\date{}

\begin{document}
\maketitle
\thispagestyle{fancy}

\section{Numerical Methods}
\subsection{Newton's Method for Root Finding}
Find $x$ such that $f(x) = 0$ using the iteration:
\[
x_{n+1} = x_n - \frac{f(x_n)}{f'(x_n)}
\]

\begin{lstlisting}
#include <stdio.h>
#include <math.h>

double f(double x)  { return x * x - 2.0; }
double fp(double x) { return 2.0 * x; }

int main(void) {
    double x = 1.0;
    for (int i = 0; i < 10; i++) {
        x = x - f(x) / fp(x);
    }
    printf("sqrt(2) ~ %.10f\n", x);
    return 0;
}
\end{lstlisting}

\section{Matrix Operations}
\begin{lstlisting}
#include <stdio.h>

#define N 3

void mat_mul(double A[N][N], double B[N][N],
             double C[N][N]) {
    for (int i = 0; i < N; i++)
        for (int j = 0; j < N; j++) {
            C[i][j] = 0.0;
            for (int k = 0; k < N; k++)
                C[i][j] += A[i][k] * B[k][j];
        }
}
\end{lstlisting}

\section{Molecular Dynamics Simulations}
Simulate particle interactions using Newton's equations of motion. The Verlet
integration scheme updates positions and velocities:
\[
x(t + \Delta t) = 2x(t) - x(t - \Delta t) + a(t)\,\Delta t^2
\]

\section{Physics Simulations}
Euler's method for a simple projectile:
\begin{lstlisting}
#include <stdio.h>

int main(void) {
    double x = 0, y = 0;
    double vx = 10.0, vy = 20.0;
    double g = 9.81, dt = 0.01;
    while (y >= 0.0) {
        x += vx * dt;
        y += vy * dt;
        vy -= g * dt;
    }
    printf("Range: %.2f m\n", x);
    return 0;
}
\end{lstlisting}

\section{Performance Optimization}
\begin{itemize}
  \item Use cache-friendly memory access patterns (row-major traversal).
  \item Minimize branch mispredictions in inner loops.
  \item Compile with \texttt{-O2} or \texttt{-O3} for optimizations.
  \item Profile with \texttt{gprof} or \texttt{perf} before optimizing.
\end{itemize}

\end{document}
